\documentclass[11pt,a4paper]{article}
\usepackage{fontspec,amsmath,amssymb,amsthm,microtype}
\setmainfont{lmroman10-regular.otf}[BoldFont=lmroman10-bold.otf,ItalicFont=lmroman10-italic.otf,BoldItalicFont=lmroman10-bolditalic.otf]
\usepackage[margin=25mm]{geometry}
\usepackage[unicode,hidelinks]{hyperref}
\setlength{\emergencystretch}{3em}
\setlength{\parskip}{0.4em}
\setlength{\parindent}{0pt}
\allowdisplaybreaks[2]
\raggedbottom
\title{\texorpdfstring{Measurable transitions and iterated particle laws}{Measurable transitions and iterated particle laws}}
\author{Omer Deniz Akyildiz\\[0.4em]\normalsize Prepared with AI assistance}
\date{}
\begin{document}\maketitle
Draft mathematical exposition; these notes may contain mistakes.


The pathwise algorithm maps become probability transitions by integrating them against independent fresh-noise measures. This section reuses the generic Markov-law construction and independence argument already proved for ULA. For KIPLMC1 it accepts a specified correlated phase innovation law; KIPLMC2 uses an explicit pair of independent standard Gaussian vectors for its two OU half steps. Law equality is kept distinct from convergence rates.



Both explicit particle recursions are identified with their independent-noise Markov iterates. The exponential transition retains correlated response pairs, while OBABO retains two distinct independent thermostat noises.


\section{\texorpdfstring{Construct the OBABO noise law and prove both step maps measurable}{Construct the OBABO noise law and prove both step maps measurable}}


the two independent thermostat noises is the product of two standard Gaussian measures on the particle Euclidean space. Its product probability property gives total mass one. Lipschitzness of the particle force implies continuity. Unfolding the exponential or OBABO map shows that each is built from this continuous force, coordinate projections, addition and scalar multiplication, so each uncurried state–noise map is continuous and hence measurable. The argument supplies force continuity first, applies measurability of continuous maps, and discharges the finite composition with continuity and measurability.



The two thermostat noises form an independent product of Gaussian vector measures. Continuity of the force makes each finite-stage update measurable. Both particle algorithms therefore define probability transitions by pushing forward the state/\allowbreak{}noise product.


\section{\texorpdfstring{Define numerical laws by independent-\allowbreak{}noise pushforwards}{Define numerical laws by independent-noise pushforwards}}


the KIPLMC1 laws applies the shared Markov recursion with the exponential step, an initial phase measure \(\mu\), and a supplied innovation measure \(\eta\). the KIPLMC2 laws uses the OBABO step and its fixed two-Gaussian noise measure. Each successor law is the pushforward of the product of the current-state law and the fresh-noise law. The probability theorems invoke the proved step measurability and probability of the initial and innovation measures, so every iterate is a probability measure. For KIPLMC1 the covariance specification of \(\eta\) remains explicit and separate; probability alone does not identify the intended correlated Gaussian law.



At each step take a fresh innovation independent of the present state and apply the specified map. Integrating this update against its innovation measure defines one transition. Repeated pushforward gives the numerical laws without assuming an underlying chain with the desired estimates.


\section{\texorpdfstring{Prove the explicit exponential recursion has the Markov law}{Prove the explicit exponential recursion has the Markov law}}


First show that the KIPLMC1 sequence equals the generic noise-driven recursion. Induction is exact: both start at the same phase and their successor expressions use the same step and innovation number. The law theorem then assumes a measurable initial random phase, a jointly independent innovation sequence with common law \(\eta\), and independence of initialization from the whole sequence. The shared generic result derives the law at every step; no additional per-step state–noise independence premise is inserted. Rewriting the explicit recursion by the proved equality yields the pushforward identity with the KIPLMC1 laws. This theorem establishes actual-law correspondence, while a quantitative sampling estimate still requires moment and coupling analysis for the specified Gaussian innovation.



The explicit recursion is a measurable function of initialization and the innovation prefix. Independence of the next block gives the one-step transition identity. Induction then identifies its law with the iterated exponential transition measure.


\section{\texorpdfstring{The independent-\allowbreak{}Gaussian OBABO trajectory has the constructed transition law}{The independent-Gaussian OBABO trajectory has the constructed transition law}}


the KIPLMC2 sequence iterates the OBABO map with the two thermostat vectors from each supplied noise pair. If these pairs are independent across time, have the product-standard-Gaussian law the two independent thermostat noises, and are independent of initialization, then the pushforward of the pathwise recursion equals the KIPLMC2 laws at every step. The generic independence theorem derives independence of the current state from the fresh pair from these global assumptions. The proof is a direct reuse of that theorem with the established measurable OBABO map. Independence within each pair is encoded by its product Gaussian law; it is not confused with the independent sequence of pairs. This completes the algorithm-law correspondence, while contraction and the discrete invariant-measure bias remain separate quantitative obligations.



The same argument applies to the two-thermostat recursion, keeping the two noises distinct and independent. Their common block law is used at every step. Thus the trajectory’s pushforward agrees with the constructed OBABO Markov law.


\section{\texorpdfstring{Choose the constructed kinetic Gaussian for Algorithm 1}{Choose the constructed kinetic Gaussian for Algorithm 1}}


The generic law identity accepts a specified iid innovation measure. the Gaussian KIPLMC1 laws supplies the correlated Gaussian measure constructed in the explicit correlated particle-noise construction, whose spatial pairs have the shared kinetic Gram covariance. Thus the transition and its particle scaling are explicit. The earlier iid realization theorem specializes to this measure, with ordinary independence of initialization and fresh innovations. A convergence theorem must still derive a target comparison for these laws; the definition does not encode one.



For Algorithm 1 the innovation law is the constructed exponential response Gaussian, including its off-diagonal covariance. Substituting that law specializes the generic transition to the correlated-Gaussian sampler. Independent velocity and position noises would define a different law.


\end{document}
